\section{Deferred Proofs}
\label{appx:proof}
\label{appx:auxiliary-proofs}

We provide the completeness, soundness, and efficiency proofs
deferred from Sections~\ref{sec:blocks}, \ref{sec:pcs}, and
\ref{sec:together}. Each proof uses the notation and parameters
of its corresponding statement.

\subsection{Proof of Proposition~\ref{proposition:ip-complexity}}
\label{appx:proof-ip}

\begin{proof}
Expanding Eq.~\eqref{eq:ip-laurent-polynomial} gives
\[
    f(X)
    =
    \sum_{\tau=-(\kappa-1)}^{\kappa-1}
    f_\tau X^\tau,
\]
where
\[
    f_\tau
    =
    \sum_{\substack{j,j'\in[\kappa]\\j-j'=\tau}}
    \left\langle
        \bm a^{(j)},\bm b^{(j')}
    \right\rangle.
\]
Its constant coefficient is
\[
    f_0
    =
    \sum_{j=1}^{\kappa}
    \left\langle
        \bm a^{(j)},\bm b^{(j)}
    \right\rangle
    =
    \langle\bm a,\bm b\rangle.
\]
Define
\[
    h_+(X)
    \coloneqq
    \sum_{\tau=1}^{\kappa-1}
    f_\tau X^{\tau-1},
    \qquad
    h_-(X)
    \coloneqq
    \sum_{\tau=1}^{\kappa-1}
    f_{-\tau}X^{\tau-1}.
\]
Then
\[
    f(X)
    =
    f_0
    +Xh_+(X)
    +X^{-1}h_-(X^{-1}),
\]
which proves completeness.

For soundness, suppose that
\(c\neq\langle\bm a,\bm b\rangle\). Because \(h_+\) and \(h_-\)
are fixed before \(\ext{x}\) is sampled,
\[
    f(X)-c-Xh_+(X)-X^{-1}h_-(X^{-1})
\]
is a nonzero Laurent polynomial whose constant coefficient is
\(f_0-c\neq0\). Multiplying it by \(X^{\kappa-1}\) produces a
nonzero polynomial of degree at most \(2\kappa-2\). Hence, the
random evaluation in Eq.~\eqref{eq:packed-inner-product-check}
accepts with probability at most
\[
    \frac{2\kappa-2}{|\ExtensionField|-1}.
\]
By Fact~\ref{fact:ip}, the subsequent length-\(m\)
inner-product reduction has soundness error at most
\(2\log_2m/|\ExtensionField|\). The claimed bound follows from
the union bound.

It remains to count the prover operations. Computing \(h_+\) and
\(h_-\) requires one length-\(m\) base-field inner product for
each ordered pair \(j\neq j'\). There are
\(\kappa(\kappa-1)\) such pairs, giving exactly
\[
    m\kappa(\kappa-1)
    =
    N(\kappa-1)
\]
multiplications of type \(\bb\).

Computing each packed vector evaluates \(\kappa\) base-field
blocks at an extension-field point and therefore requires
\(m\kappa=N\) multiplications of type \(\be\). The two packed
vectors consequently require \(2N\,\be\).

By Fact~\ref{fact:ip}, reducing the resulting length-\(m\)
extension-field inner product requires
\[
    4m+O(\log m)
\]
multiplications of type \(\ee\). Computing the positive and
negative powers of \(\ext{x}\) contributes \(O(\kappa)\)
additional \(\ee\) operations. Substituting \(m=N/\kappa\) gives
\[
    N(\kappa-1)\,\bb
    +
    2N\,\be
    +
    \left(
        \frac{4N}{\kappa}
        +
        O\left(
            \kappa+\log\frac{N}{\kappa}
        \right)
    \right)\ee.
\]

Finally, \(\bm A\) and \(\bm B\) occupy \(O(N)\) base-field
coordinates. The cross-block inner products can be streamed while
retaining only the \(2\kappa-2\) coefficients of \(h_+\) and
\(h_-\). The challenge powers occupy \(O(\kappa^2)\) base-field
coordinates, and the two packed vectors occupy \(O(N)\).
Using in-place inner-product reduction, the peak prover memory is
therefore \(O(N+\kappa^2)\), which is \(O(N)\) when
\(m\geq\kappa\).


\emph{Soundness.}
Because \(h_+\) and \(h_-\) are fixed before \(\overline{x}\) is
sampled, an incorrect claim \(c\neq\langle\bm a,\bm b\rangle\)
makes
\[
    f(X)-c-Xh_+(X)-X^{-1}h_-(X^{-1})
\]
a nonzero Laurent polynomial: its constant coefficient is
\(f_0-c\neq0\). Multiplying by \(X^{\kappa-1}\) produces a
nonzero polynomial of degree at most \(2\kappa-2\). Hence,
Eq.~\eqref{eq:packed-inner-product-check} holds for a uniformly
random \(\overline{x}\in\ExtensionField^{\times}\) with
probability at most
\[
    \frac{2\kappa-2}{|\ExtensionField|-1}.
\]
The total soundness error additionally includes that of the
underlying extension-field inner-product proof.
\end{proof}

\subsection{Proof of Proposition~\ref{proposition:hadamard-complexity}}
\label{appx:proof-hadamard}

\begin{proof}
For completeness, define
\[
    \ext{d}_j
    \coloneqq
    \left\langle
        \ext{\bm l},
        \bm a^{(j)}\circ\bm b^{(j)}
    \right\rangle
    \qquad
    \text{for }j\in[\kappa].
\]
For an honest prover,
\(\ext{\bm c}=\ext{\bm d}\), and therefore
\[
    \ext{c}
    =
    \left\langle
        \ext{\bm c},\ext{\bm r}
    \right\rangle
    =
    \ext{\bm l}^{\top}
    (\bm A\circ\bm B)
    \ext{\bm r}.
\]
Moreover, the residual polynomial satisfies
\[
    f(\alpha_j)
    =
    \ext{d}_j-\ext{c}_j
    =
    0
    \qquad
    \text{for every }j\in[\kappa].
\]
Thus, \(\Psi(X)\) divides \(f(X)\). Since
\[
    \deg f\leq2\kappa-2
    \qquad\text{and}\qquad
    \deg\Psi=\kappa,
\]
the quotient
\[
    h(X)=\frac{f(X)}{\Psi(X)}
\]
has degree at most \(\kappa-2\).

Evaluating the quotient identity at \(\ext{x}\) gives
\[
\begin{aligned}
    \ext{d}
    &=
    g_C(\ext{x})
    +
    \Psi(\ext{x})h(\ext{x})\\
    &=
    \left\langle
        \ext{\bm a},
        \ext{\bm b}\circ\ext{\bm l}
    \right\rangle.
\end{aligned}
\]
Hence, the claim passed to the inner-product reduction is correct.

Let
\[
    \ext{\bm\chi}_+
    \coloneqq
    \tensorspan{\ext{\bm y}'},
    \qquad
    \ext{\bm\chi}_-
    \coloneqq
    \tensorspan{(\ext{\bm y}')^{-1}}.
\]
The terminal scalars returned by Fact~\ref{fact:ip} satisfy
\[
\begin{aligned}
    \ext{u}_1
    &=
    \left\langle
        \ext{\bm a},\ext{\bm\chi}_+
    \right\rangle=
    \sum_{j=1}^{\kappa}
    \psi_j(\ext{x})
    \left\langle
        \bm a^{(j)},\ext{\bm\chi}_+
    \right\rangle
    =
    \left\langle
        \ext{\bm v}_a,\ext{\bm\rho}
    \right\rangle,
\end{aligned}
\]
and
\[
\begin{aligned}
    \ext{u}_2
    &=
    \left\langle
        \ext{\bm b}\circ\ext{\bm l},
        \ext{\bm\chi}_-
    \right\rangle
    =
    \sum_{j=1}^{\kappa}
    \psi_j(\ext{x})
    \left\langle
        \bm b^{(j)},
        \ext{\bm l}\circ\ext{\bm\chi}_-
    \right\rangle
    =
    \left\langle
        \ext{\bm v}_b,\ext{\bm\rho}
    \right\rangle.
\end{aligned}
\]
Therefore, the Lagrange consistency checks accept.

Let
\[
    \ext{\bm\chi}''
    \coloneqq
    \tensorspan{\ext{\bm y}''}.
\]
The randomized column aggregation gives
\[
\begin{aligned}
    \ext{v}_1
    &=
    \left\langle
        \ext{\bm v}_a,\ext{\bm\chi}''
    \right\rangle
    =
    \left\langle
        \bm a,
        \ext{\bm\chi}''
        \kronecker
        \ext{\bm\chi}_+
    \right\rangle
    =
    \left\langle
        \bm a,
        \tensorspan{\ext{\bm y}_a}
    \right\rangle.
\end{aligned}
\]
Similarly, since
\[
    \ext{\bm l}
    =
    \tensorspan{\ext{\bm z}_l},
\]
we have
\[
    \ext{\bm l}\circ\ext{\bm\chi}_-
    =
    \tensorspan{
        \ext{\bm z}_l
        \circ
        (\ext{\bm y}')^{-1}
    }.
\]
It follows that
\[
\begin{aligned}
    \ext{v}_2
    &=
    \left\langle
        \ext{\bm v}_b,\ext{\bm\chi}''
    \right\rangle
    =
    \left\langle
        \bm b,
        \ext{\bm\chi}''
        \kronecker
        \left(
            \ext{\bm l}\circ\ext{\bm\chi}_-
        \right)
    \right\rangle
    =
    \left\langle
        \bm b,
        \tensorspan{\ext{\bm y}_b}
    \right\rangle,
\end{aligned}
\]
which proves completeness of the terminal MLE claims.

For soundness, suppose that the weighted Hadamard-product claim is
false. If
\[
    \ext{c}
    =
    \left\langle
        \ext{\bm c},\ext{\bm r}
    \right\rangle,
\]
then \(\ext{\bm c}\neq\ext{\bm d}\), and hence
\(f(\alpha_j)\neq0\) for at least one \(j\in[\kappa]\). Therefore, for
any \(h(X)\) fixed before \(\ext{x}\) is sampled,
\[
    f(X)-\Psi(X)h(X)
\]
is a nonzero polynomial of degree at most \(2\kappa-2\).
Consequently, a uniformly random
\(\ext{x}\in\ExtensionField\setminus\mathcal H\) satisfies the
quotient identity with probability at most
\[
    \frac{2\kappa-2}
         {|\ExtensionField|-\kappa}.
\]

By Fact~\ref{fact:ip}, the subsequent length-\(m\)
inner-product reduction contributes soundness error at most
\[
    \frac{2\log_2m}
         {|\ExtensionField|}.
\]

It remains to consider the column values sent after the
inner-product reduction. Let
\[
\begin{aligned}
    \widehat v_{a,j}
    &\coloneqq
    \left\langle
        \bm a^{(j)},\ext{\bm\chi}_+
    \right\rangle,
    \qquad
    \widehat v_{b,j}
    \coloneqq
    \left\langle
        \bm b^{(j)},
        \ext{\bm l}\circ\ext{\bm\chi}_-
    \right\rangle
\end{aligned}
\]
denote their correct values. If either
\(\ext{\bm v}_a\neq\widehat{\bm v}_a\) or
\(\ext{\bm v}_b\neq\widehat{\bm v}_b\), the corresponding
difference vector defines a nonzero multilinear polynomial in
\(\log_2\kappa\) variables. Since these vectors are fixed before
\(\ext{\bm y}''\) is sampled, the randomized MLE aggregation
conceals an incorrect vector with probability at most
\[
    \frac{\log_2\kappa}{|\ExtensionField|}
\]
for each terminal claim. A union bound over the two claims gives
the additional error
\[
    \frac{2\log_2\kappa}{|\ExtensionField|}.
\]
Combining the three error terms proves the stated soundness bound.

We finally count the prover operations. Computing
\(\ext{\bm c}\) requires \(N\) multiplications of type \(\bb\)
to form \(\bm A\circ\bm B\), followed by \(N\) multiplications
of type \(\be\) to contract the result with
\(\ext{\bm l}\).

For each of the \(m\) rows, interpolating
\(g_{A,i}(X)\) and \(g_{B,i}(X)\) requires
\(2\kappa^2\) multiplications of type \(\bb\), while multiplying
the resulting degree-\((\kappa-1)\) polynomials requires another
\(\kappa^2\) multiplications of type \(\bb\). Thus, the total
base-field cost is
\[
    N+3m\kappa^2
    =
    (3\kappa+1)N
\]
multiplications of type \(\bb\).

Each product
\(g_{A,i}(X)g_{B,i}(X)\) contains \(2\kappa-1\)
coefficients. Contracting these coefficients with
\(\ext{l}_i\) over all \(m\) rows requires
\[
    m(2\kappa-1)
    =
    2N-m
\]
multiplications of type \(\be\). Interpolating \(g_C(X)\) and
dividing the residual polynomial by \(\Psi(X)\) require
\(O(\kappa^2)\) additional \(\be\) operations.

After \(\ext{x}\) is sampled, forming
\(\ext{\bm a}=\bm A\ext{\bm\rho}\) and
\(\ext{\bm b}=\bm B\ext{\bm\rho}\) requires \(2N\,\be\).
After the inner-product reduction, computing the two vectors of
column values \(\ext{\bm v}_a\) and \(\ext{\bm v}_b\) requires
another \(2N\,\be\). Including the \(N\,\be\) operations used to
construct \(\ext{\bm c}\), the cross-field cost is therefore
\[
\begin{aligned}
    &N+(2N-m)+2N+2N\\
    &\qquad=
    7N-m
    =
    7N-\frac{N}{\kappa}
\end{aligned}
\]
multiplications of type \(\be\), plus
\(O(\kappa^2)\) lower-order operations.

Forming
\(\ext{\bm b}\circ\ext{\bm l}\) before the inner-product
reduction requires \(m\,\ee\). By Fact~\ref{fact:ip}, the
length-\(m\) reduction requires
\[
    4m+O(\log m)
\]
additional \(\ee\) operations. Computing
\[
    \ext{\bm l}
    \circ
    \tensorspan{(\ext{\bm y}')^{-1}}
\]
for the column claims requires another \(m\,\ee\).
Evaluating the Lagrange basis and the degree-\(O(\kappa)\)
polynomials, and forming the random column aggregations, requires
\(O(\kappa)\) further \(\ee\) operations. Hence, the total
extension-field cost is
\[
    6m+O(\kappa+\log m)
    =
    \frac{6N}{\kappa}
    +
    O\left(
        \kappa+\log\frac{N}{\kappa}
    \right)
\]
multiplications of type \(\ee\).

Finally, the matrices \(\bm A\) and \(\bm B\) occupy \(2N\)
base-field coordinates. The prover processes the row polynomials
sequentially and stores only \(O(\kappa)\) extension-field
coefficients for the residual and quotient polynomials. The
column-value vectors and random tensor weights require
\(O(\kappa)\) extension-field elements, while the packed vectors
and inner-product reduction require a constant number of
length-\(m\) extension-field buffers. The peak prover memory is
therefore
\[
    O(N+\kappa^2)
\]
base-field coordinates, which is \(O(N)\) whenever
\(m\geq\kappa\).
\end{proof}

\subsection{Proof of Proposition~\ref{proposition:permutation-complexity}}
\label{appx:proof-permutation}

\begin{proof}
For completeness, if \(\bm a\) is a permutation of \(\bm b\),
then
\[
    \prod_{i=1}^{N}(\ext z+a_i)
    =
    \prod_{i=1}^{N}(\ext z+b_i)
\]
for every \(\ext z\in\ExtensionField\). It therefore suffices to
prove completeness of one grand-product reduction.

By construction of the column-product tree,
\[
    \bm C_{t+1}(\ext z)
    =
    \bm C_{t,L}(\ext z)
    \circ
    \bm C_{t,R}(\ext z)
\]
for every
\(t=0,\ldots,\log_2\kappa-1\). The recursive row-product tree
similarly satisfies
\[
    \ext{\bm c}_{i+1}
    =
    \ext{\bm c}_{i,L}
    \circ
    \ext{\bm c}_{i,R},
    \qquad
    i=0,\ldots,\log_2m-1,
\]
and hence
\[
    \ext c_{\log_2m}
    =
    \mathsf{GrandProd}
    \bigl(
        \bm C_{\log_2\kappa}(\ext z)
    \bigr).
\]

Every weighted Hadamard claim in these two trees is proved
directly using the inner-product proof. For
\[
    \ext{\bm u},
    \ext{\bm v},
    \ext{\bm\xi}
    \in\ExtensionField^n,
\]
we rewrite
\begin{equation}
\begin{aligned}
    \ext c
    &\EqOrNot
    \left\langle
        \ext{\bm\xi},
        \ext{\bm u}\circ\ext{\bm v}
    \right\rangle\\
    &=
    \left\langle
        \ext{\bm u},
        \ext{\bm v}\circ\ext{\bm\xi}
    \right\rangle.
\end{aligned}
\label{eq:permutation-hadamard-as-ip}
\end{equation}
The prover forms
\(
    \ext{\bm v}\circ\ext{\bm\xi}
\)
and applies the inner-product reduction of
Fact~\ref{fact:ip} to
Eq.~\eqref{eq:permutation-hadamard-as-ip}. Thus, the Hadamard
relations require no separate proof primitive.

Applying this reduction recursively first unfolds the
row-product tree and then the column-product tree. The final
claim is
\[
    \left\langle
        \ext{\bm\xi}_0,\bm d
    \right\rangle
    \EqOrNot
    \ext c_0
    -
    \ext z
    \left\langle
        \ext{\bm\xi}_0,\bm 1_N
    \right\rangle,
\]
which is an MLE opening of the original committed vector
\(\bm d\). This proves completeness.

For soundness, suppose first that
\(\bm a\) is not a permutation of \(\bm b\). Then
\[
    Q(Z)
    \coloneqq
    \prod_{i=1}^{N}(Z+a_i)
    -
    \prod_{i=1}^{N}(Z+b_i)
\]
is a nonzero polynomial. Its degree is at most \(N-1\), since
the two leading degree-\(N\) terms cancel. Therefore, the
fingerprint check holds at a uniformly random
\(\ext z\in\ExtensionField\) with probability at most
\[
    \min\left\{
        1,
        \frac{N-1}{|\ExtensionField|}
    \right\}.
\]

We next bound the soundness error of one grand-product proof.
At row level \(i\), the inner-product relation in
Eq.~\eqref{eq:permutation-hadamard-as-ip} has length
\[
    \frac{m}{2^{i+1}}.
\]
By Fact~\ref{fact:ip}, its soundness error is at most
\[
    \frac{
        2\log_2(m/2^{i+1})
    }{
        |\ExtensionField|
    }.
\]
Summing over the row-product tree gives
\[
\begin{aligned}
    \sum_{i=0}^{\log_2m-1}
    \frac{
        2\log_2(m/2^{i+1})
    }{
        |\ExtensionField|
    }
    &=
    \frac{
        \log_2m(\log_2m-1)
    }{
        |\ExtensionField|
    }.
\end{aligned}
\]

At column level \(t\), the interpolation domain contains
\[
    2^{\log_2\kappa-t}
\]
points, and the corresponding residual polynomial has degree at
most
\[
    2^{\log_2\kappa-t+1}-2.
\]
Because the quotient polynomial is fixed before the evaluation
challenge is sampled, the polynomial check fails with
probability at most
\[
    \frac{
        2^{\log_2\kappa-t+1}-2
    }{
        |\ExtensionField|
        -
        2^{\log_2\kappa-t}
    }.
\]

The evaluated Hadamard claim at this level has length
\[
    \frac{m}{2^{\log_2\kappa-t}}.
\]
The inner-product reduction of Fact~\ref{fact:ip} therefore
contributes
\[
    \frac{
        2\log_2
        \left(
            m/2^{\log_2\kappa-t}
        \right)
    }{
        |\ExtensionField|
    }.
\]
The random combination of the resulting column claims
contributes at most
\[
    \frac{
        2\log_2
        \left(
            2^{\log_2\kappa-t}
        \right)
    }{
        |\ExtensionField|
    }.
\]
Their sum is
\[
    \frac{2\log_2m}{|\ExtensionField|}.
\]
Summing over
\(t=0,\ldots,\log_2\kappa-1\) gives the stated bound on
\(\epsilon_{\mathsf{GP}}\). The full soundness bound follows by
applying the grand-product proof to both committed vectors and
using the union bound.

We finally analyze the prover complexity. At column level \(t\),
constructing the coefficient representation and its quotient
polynomial requires
\[
    O\left(
        \frac{
            \kappa^2m
        }{
            2^{\log_2\kappa-t}
        }
    \right)\bb.
\]
Consequently,
\[
\begin{aligned}
    \sum_{t=0}^{\log_2\kappa-1}
    O\left(
        \frac{
            \kappa^2m
        }{
            2^{\log_2\kappa-t}
        }
    \right)
    &=
    O\bigl(m\kappa(\kappa-1)\bigr)\\
    &=
    O(N\kappa)
\end{aligned}
\]
base-field multiplications. Evaluating the coefficient
polynomials and constructing the packed vectors require
\(O(N)\) multiplications of type \(\be\).

For a length-\(n\) weighted Hadamard claim, forming
\(
    \ext{\bm v}\circ\ext{\bm\xi}
\)
requires \(n\) multiplications of type \(\ee\). By
Fact~\ref{fact:ip}, the subsequent inner-product reduction
requires
\[
    4n+O(\log n)
\]
additional \(\ee\) multiplications. Hence, directly proving a
length-\(n\) Hadamard claim costs
\[
    5n+O(\log n)
\]
multiplications of type \(\ee\).

The lengths of the row-tree inner products sum to
\[
    \sum_{i=0}^{\log_2m-1}
    \frac{m}{2^{i+1}}
    =
    m-1.
\]
Thus, the row-tree reductions require
\[
    5(m-1)
    +
    O\bigl((\log_2m)^2\bigr)
\]
multiplications of type \(\ee\).

The active inner-product lengths in the column reductions sum to
\[
\begin{aligned}
    \sum_{t=0}^{\log_2\kappa-1}
    \frac{m}{2^{\log_2\kappa-t}}
    &=
    m\frac{\kappa-1}{\kappa}.
\end{aligned}
\]
Their leading cost is therefore
\[
    5m\frac{\kappa-1}{\kappa}
\]
multiplications of type \(\ee\). Including the lower-order
challenge-generation and polynomial-evaluation costs gives
\[
\begin{aligned}
    &5(m-1)
    +
    5m\frac{\kappa-1}{\kappa}
    +
    O\left(
        \kappa
        +
        (\log_2m)^2
        +
        \log_2m\log_2\kappa
    \right)\\
    &\qquad=
    \frac{10N}{\kappa}
    +
    O\left(
        \kappa
        +
        (\log_2m)^2
        +
        \log_2m\log_2\kappa
    \right)
\end{aligned}
\]
multiplications of type \(\ee\).

The prover constructs and discards the coefficient layers
sequentially. The original vector occupies \(O(N)\) base-field
coordinates, while the quotient coefficients and challenge
powers occupy \(O(\kappa^2)\) base-field coordinates. Therefore,
the peak prover memory is \(O(N+\kappa^2)\), which is \(O(N)\)
when \(m\geq\kappa\).
\end{proof}

\subsection{Proof of Proposition~\ref{proposition:packed-pcs}}
\label{appx:proof-pcs}

\begin{proof}
\noindent\textbf{Binding and completeness.}
Each row of \(\bm A\) is the unique base-field coordinate vector
of the corresponding entry of \(\bm A\ext{\bm\beta}\).
Thus two distinct witness matrices cannot produce the same packed
vector, and commitment binding follows from the underlying PCS.
Commitment hiding is inherited by applying the underlying commitment
to this fixed encoding.

For an honest prover,
\(\bm V=\bm A^\top\bm\Xi\) and
\(\bm W=\bm U^\top\bm\Xi\).
Eq.~\eqref{eq:pcs-encoded-matrix} gives the extension-field
inner-product relation, and Eq.~\eqref{eq:pcs-query-matrix} gives
its query-side terminal value.
Every base-field projection in Eq.~\eqref{eq:pcs-base-projection}
is then correct.
Eq.~\eqref{eq:pcs-public-query-evaluation} recovers the exact
base-field rows needed in Eq.~\eqref{eq:pcs-implicit-query-check},
so all terminal checks accept.
The remaining MLE claim is an honest opening of
\(\ext{\bm a}\), establishing completeness.

\noindent\textbf{Soundness of the extension-field reduction.}
Fix the matrix \(\bm A\) bound by the commitment and suppose that
Eq.~\eqref{eq:pcs-opening-relation} is false.
If the first scalar check accepts, then
\(\ext{\bm v}\neq\bm A^\top\ext{\bm\chi}\), hence
\(\bm V\neq\bm A^\top\bm\Xi\).
By injectivity of basis packing, the following row is nonzero:
\[
    \ext{\bm\beta}^{\top}
    (\bm V-\bm A^\top\bm\Xi)
    \in\ExtensionField^{1\times\kappa}\setminus\{\bm0_\kappa^\top\}.
\]
It is fixed before the uniform challenge \(\ext{\bm r}\).
Its inner product with \(\ext{\bm r}\) vanishes with probability
\(1/|\ExtensionField|\).
Outside this event, Eq.~\eqref{eq:pcs-random-opening} is false.

Fact~\ref{fact:ip} contributes error
\(2\log_2m/|\ExtensionField|\).
Outside this reduction error, an accepting inner-product transcript
leaves at least one false terminal claim in
Eq.~\eqref{eq:pcs-terminal-claims}.
A false witness-side claim is accepted by the underlying PCS with
probability at most \(\epsilon_{\mathsf{PCS}}\).
It remains to bound acceptance when the query-side claim is false.

\noindent\textbf{Soundness of the base-field query check.}
The matrix \(\bm U\), the coordinate matrix \(\bm\Xi\), and the
claimed scalar \(\ext w\) are fixed by the preceding transcript.
If the scalar check in Eq.~\eqref{eq:pcs-query-matrix} accepts a
false query-side value, then
\(\bm W\neq\bm U^\top\bm\Xi\).
Let
\(\bm D=\bm W-\bm U^\top\bm\Xi\neq\bm0_{\kappa\times\kappa}\).
For fresh independent uniform
\(\bm s_{(i)},\bm t_{(i)}\in\Field^\kappa\), the two-sided
projection fails to detect this discrepancy with probability at most:
\[
\begin{aligned}
    \Pr[\bm s_{(i)}^\top\bm D\bm t_{(i)}=0]
    &\leq
       \Pr[\bm D\bm t_{(i)}=\bm0_\kappa]
       +\frac{1}{|\Field|}\\
    &\leq\frac{2}{|\Field|}.
\end{aligned}
\]
Indeed, a nonzero row of \(\bm D\) vanishes on a uniform
\(\bm t_{(i)}\) with probability \(1/|\Field|\).
Conditioned on \(\bm D\bm t_{(i)}\neq\bm0_\kappa\), its inner
product with the independent \(\bm s_{(i)}\) vanishes with the
same probability.

Outside the projection event,
Eq.~\eqref{eq:pcs-base-projection} is a false inner-product claim.
Applying Fact~\ref{fact:ip} over \(\Field\) contributes at most
\(2\log_2m/|\Field|\).
Its terminal claims are checked exactly from the public queries by
Eq.~\eqref{eq:pcs-implicit-query-check}.
One repetition therefore accepts with probability at most
\((2\log_2m+2)/|\Field|\).

The same nonzero \(\bm D\) is used throughout the repetitions.
Conditioned on any earlier transcript, the next repetition has fresh
projection and inner-product challenges and the same error bound,
even if the prover's responses depend on the earlier transcript.
Thus all \(k\) repetitions accept with probability at most:
\[
    \left(\frac{2\log_2m+2}{|\Field|}\right)^k
    \leq2^{-\lambda}.
\]
Combining the extension-field projection, its inner-product reduction,
the terminal PCS opening, and this query-side check gives
Eq.~\eqref{eq:pcs-soundness}.

\noindent\textbf{Prover arithmetic and memory.}
The prover expands the public tensor queries
\(\ext{\bm\chi}\) and \(\ext{\bm\xi}\), and reads their basis
coordinates to obtain \(\bm\Xi\) and \(\bm U\).
This costs \(O(m)\) extension-field operations.
Computing each of the two matrix products
\(\bm V=\bm A^\top\bm\Xi\) and
\(\bm W=\bm U^\top\bm\Xi\) takes \(\kappa N\) base-field
multiplications.
The message \(\ext{\bm v}=\bm V\ext{\bm\beta}\) is obtained by
packing the rows of \(\bm V\).

Computing \(\bm\Xi\ext{\bm r}\) takes \(N\) mixed-field
multiplications; the extension-field inner-product reduction takes
\(4m+O(\log m)\) extension-field multiplications.
In each repetition, forming \(\bm U\bm s_{(i)}\) and
\(\bm\Xi\bm t_{(i)}\) takes \(2N\) base-field multiplications,
and the base-field inner-product reduction takes
\(4m+O(\log m)\) more.
Computing the public scalar target
\(\bm s_{(i)}^\top\bm W\bm t_{(i)}\) adds \(O(\kappa^2)\)
base-field operations per repetition.
The extension-field scalar targets can be evaluated by first
packing the columns of \(\bm V\) or \(\bm W\), at a total cost
of \(O(\kappa)\) extension-field multiplications.
These costs give Eq.~\eqref{eq:pcs-prover-cost}.

The matrices \(\bm A,\bm\Xi,\bm U\) occupy \(O(N)\) base-field
coordinates, and \(\bm V,\bm W\) occupy \(O(\kappa^2)\).
A constant number of length-\(m\) extension-field buffers also
occupy \(O(m\kappa)=O(N)\) base-field coordinates.
Generating the base-field operands and processing the repetitions
sequentially gives the stated peak memory bound.

\noindent\textbf{Verifier arithmetic and communication.}
The verifier never expands \(\bm U\) or \(\bm\Xi\).
The initial scalar check and the two extension-field projection
targets require \(O(\kappa)\) extension-field multiplications,
after packing the columns of \(\bm V\) and \(\bm W\).
The outer inner-product verifier uses \(O(\log m)\)
extension-field operations.

Each base-field repetition requires \(O(\kappa^2)\) base-field
operations for its scalar target, \(O(\log m)\) for the
inner-product verifier, and \(O(\kappa)\) for the two terminal checks.
The two tensor contractions use \(O(\log m)\) mixed-field and
extension-field operations.
This gives Eq.~\eqref{eq:pcs-verifier-cost}.

The prover sends \(\ext{\bm v}\), consisting of \(\kappa\)
extension-field elements, and \(\bm W\), consisting of
\(\kappa^2\) base-field elements.
The outer inner-product transcript adds \(2\log_2m+2\)
extension-field elements; the \(k\) inner transcripts add
\(k(2\log_2m+2)\) base-field elements.
There is no separate transmission of \(\bm V\), since it is the
coordinate representation of \(\ext{\bm v}\).
This proves the communication bound.
\end{proof}

\subsection{Proof of Proposition~\ref{proposition:position-encoding-evaluation}}
\label{appx:proof-position-encoding}

\begin{proof}
Let \(\bm d\in\Field^{2^s}\) list the encoded digits
\(\sum_{t=1}^s\delta_tb_t\) in lexicographic order over
\((b_1,\ldots,b_s)\in\{0,1\}^s\).
Its MLE is:
\begin{equation}
    \left\langle\bm d,\tensorspan{\ext{\bm y}}\right\rangle
       =\sum_{t=1}^{s}\delta_t\ext y_t,
       \qquad \ext{\bm y}\in\ExtensionField^s.
\label{eq:digit-vector-mle}
\end{equation}
The right-hand side is multilinear and agrees with the encoded
digit at every Boolean point, proving the identity in every
characteristic.
Each base-\(\underline p\) digit occupies one consecutive
\(s\)-bit block. Prepending \(\gamma s-\mu\) zero bits restricts
the full digit domain to positions \(0,\ldots,L-1\).
Applying Eq.~\eqref{eq:digit-vector-mle} to each block and taking
the basis-weighted sum proves
Eq.~\eqref{eq:position-encoding-evaluation}.
Only the \(\mu\) original coordinates need to be processed.

For arbitrary tensor factors
\(\ext{\bm\chi}^{(t)}=(\ext\chi^{(t)}_0,\ext\chi^{(t)}_1)^\top\),
the contraction with the \(t\)-th bit function is:
\[
    \ext\chi^{(t)}_1
       \prod_{\substack{u=1\\u\ne t}}^\mu
          (\ext\chi^{(u)}_0+\ext\chi^{(u)}_1).
\]
All such products are obtained in \(O(\mu)\) operations using
prefix and suffix products, without division.
The same basis-weighted sum therefore evaluates the position
labels in \(O(\log L)\) operations for these queries as well.
\end{proof}

\subsection{Proof of Theorem~\ref{theorem:single-instance-packing}}
\label{appx:proof-single-instance}

\begin{proof}
\noindent\textbf{Completeness and reduction correctness.}
For a satisfying assignment, the five claims in
Eq.~\eqref{eq:plonkWeightedHadamard1} and
Eq.~\eqref{eq:plonkWeightedHadamard2} sum to zero.
The vector \(\bm w_M=\bm w_L\circ\bm w_R\) is determined
by the three wire vectors, and its terminal evaluation is
exactly a weighted Hadamard-product claim on those vectors.
Completeness of the five Hadamard reductions and their terminal
openings therefore proves completeness of the gate argument.

The basis vectors \(\ext\theta_1,\ldots,\ext\theta_\gamma,
\ext\theta_d\) are linearly independent over \(\Field\).
Consequently, equality between two encoded entries forces both
their wire values and position digits to agree.
In particular, the entry of \(\ext{\bm a}\) with label
\(\ext\tau_i\) can match only the entry of \(\ext{\bm b}\)
at position \(\sigma^{-1}(i)\).
Equality of the two multisets is thus equivalent to
\(w_i=w_{\sigma^{-1}(i)}\) for every \(i\), proving
Eq.~\eqref{eq:together-copy-equivalence}.
Padding adds fixed points and does not change the relation.
Completeness of the permutation and opening protocols completes
the argument.

\noindent\textbf{Soundness.}
Condition on binding of the three witness commitments.
Since fixed-basis packing is injective, these commitments determine
three base-field wire vectors and hence the virtual \(\bm w_M\)
before any challenge is sampled.
The selectors and permutation digits are fixed by circuit setup.

If the gate residual is nonzero, its random MLE evaluation
vanishes with probability at most \(\log_2N/|\ExtensionField|\).
Otherwise, a prover satisfying the scalar sum must give at least
one false claim among the five terms.
Following its reduction either violates one of the five
Hadamard proofs or produces a false terminal opening.
This includes the nested claim on \(\bm w_M\): its value is
fixed before the nested proof samples its challenges, while
its correct value is determined by the committed wire vectors.

If a copy constraint is false, the two encoded multisets differ.
The fingerprint check in Section~\ref{sec:permutation} then has
error at most \(\min\{1,(L-1)/|\ExtensionField|\}\).
Conditioned on a nonzero fingerprint, acceptance requires failure
of one of the two grand-product reductions or a false opening.
The deterministic joint encoding contributes no additional error.

There are five selector and five wire openings in the gate
argument. Each of the two copy terminals contributes three
wire openings. Only the operand containing
\(\ext{\bm\tau}_\sigma\) contributes digit-table openings,
one for each of its \(\gamma\) base-field coordinates.
The public-I/O and identity-label evaluations are exact.
Thus there are sixteen length-\(N\) packed openings and
\(\gamma\) length-\(L\) packed openings.
Their adaptive evaluation errors are bounded by
Eq.~\eqref{eq:together-opening-error}.
A union bound over these events and witness-binding failure
gives Eq.~\eqref{eq:together-soundness}.

\noindent\textbf{Efficiency.}
The five Hadamard reductions and the sixteen length-\(N\)
packed openings cost \(O((\kappa+k)N)\bb\).
The \(\gamma\) digit-table openings cost
\(O(\gamma(\kappa+k)L)\bb\).
Both estimates follow from
Propositions~\ref{proposition:hadamard-complexity} and
\ref{proposition:packed-pcs}, since \(m\geq\kappa\) and
\(L/\kappa\geq m\).

For the copy proofs, a multiplication in \(\mathbb K\) costs
\(O(d^2)\bb\), a \(\mathbb K\)-by-\(\ExtensionField\)
multiplication costs \(O(\kappa d)\bb\), and a multiplication
in \(\ExtensionField\) costs \(O(\kappa^2)\bb\).
Substituting \(L\) and \(\kappa_{\mathbb K}=\kappa/d\)
in Proposition~\ref{proposition:permutation-complexity}
therefore gives:
\[
    O\bigl(\kappa dL+\kappa^2(\log_2L)^2\bigr)\bb.
\]
The length-dependent term absorbs the additive packing costs,
using \(L\geq N\geq\kappa^2\).
Summing these contributions proves the prover-time bound.

For each interpolation check, precomputed polynomial coefficients
allow evaluation of a complete Lagrange basis in
\(O(\kappa^3)\bb\); the tower-field analog is no larger.
The packed openings contribute
\(O((\gamma+1)\kappa^2[\kappa+k\log_2L])\bb\)
to verifier time.
The recursive grand-product checks contribute
\(O(\kappa^2(\log_2L)^2)\bb\).
The aligned wire segments require only tensor restrictions,
and the public-I/O contribution can be evaluated in
\(O(n_{\mathsf{io}}+\log L)\) extension-field operations.
Proposition~\ref{proposition:position-encoding-evaluation}
handles the identity labels. These costs give the verifier bound.

The Hadamard and grand-product reductions communicate
\(O(\kappa+(\log_2L)^2)\) extension-field elements.
Adding the packed-opening messages from
Proposition~\ref{proposition:packed-pcs} gives the communication
bound in \(\Field\)-coordinates.
Sequential execution retains \(O(N+Ld+\kappa^2)=O(N+Ld)\)
coordinates, including the digit tables and the
\(\mathbb K\)-valued copy operands.
There are three online witness commitments, and each packed
opening makes one underlying opening call at the stated length.
All underlying PCS costs and workspace are accounted for separately.
\end{proof}

\subsection{Proof of Theorem~\ref{theorem:multiple-instance-packing}}
\label{appx:proof-multiple-instance}

\begin{proof}
\noindent\textbf{Completeness and exact instance packing.}
Let \(\bm R\in\Field^{m\times\kappa}\) contain the gate
residuals of the \(\kappa\) instances as columns.
The left-hand side of Eq.~\eqref{eq:together-multiple-gates}
is exactly \(\bm R\ext{\bm\beta}\), including the term
\(\bm q_C\bm1_\kappa^\top\ext{\bm\beta}\).
Linear independence of the basis implies:
\[
    \bm R\ext{\bm\beta}=\bm0_m
       \quad\Longleftrightarrow\quad
    \bm R=\bm0_{m\times\kappa}.
\]
Similarly, packing commutes with the common row permutation,
so \(\sigma(\ext{\bm w})=\ext{\bm w}\) is equivalent to
the copy constraints in every instance.
This argument also covers the packed public I/O.
The virtual product vector is defined by
Eq.~\eqref{eq:together-multiple-wires}, and its queried value
is authenticated by Eq.~\eqref{eq:together-multiple-product-opening}.
All scalar identities and product-tree relations hold for an
honest prover, proving completeness.

\noindent\textbf{Gate soundness.}
Condition on binding of the witness commitments.
Every committed vector in \(\ExtensionField^m\) has a unique
matrix of \(\Field\)-coordinates in the fixed basis.
Thus both the gate residual and \(\ext{\bm w}_M\) are fixed
functions of the committed witnesses.
If some gate is invalid, the packed residual is nonzero.
Its random length-\(m\) tensor projection vanishes with
probability at most \(\log_2m/|\ExtensionField|\).

For a fixed projection point \(\ext{\bm x}\), the gate
sum-check polynomial is the product of the multilinear equality
weight with the residual expression formed from selector and wire
MLEs. Each selector-times-wire term has individual degree two,
and multiplication by the equality weight raises this to three.
The constant-selector term has lower degree.
The sum-check therefore contributes at most
\(3\log_2m/|\ExtensionField|\).
Its five selector and three committed-wire evaluations are
checked by direct PCS openings.
The remaining virtual-product evaluation is checked by one
packed Hadamard proof and its two packed openings.
This yields the first two algebraic terms in
Eq.~\eqref{eq:together-multiple-soundness}, without a union
bound over the \(\kappa\) instances.

\noindent\textbf{Copy fingerprint.}
Suppose that the packed assignment violates a copy constraint.
Consider the bivariate difference polynomial:
\[
    P(Y,Z)=
       \prod_{i=1}^{L}(Z+\ext w_i+Y\ext\tau_i)
       -\prod_{i=1}^{L}(Z+\ext w_i+Y\ext\tau_{\sigma,i}).
\]
The labels \(\ext\tau_i\) are distinct because the bit-to-field
digit map is injective and \(\gamma\leq\kappa\) digit slots use
independent basis coordinates.
If \(P\) were zero, unique factorization in
\(\ExtensionField[Y,Z]\) would equate the two multisets of
linear factors, each monic in \(Z\).
Matching their \(Y\)-coefficients would then force
\(\ext w_i=\ext w_{\sigma^{-1}(i)}\) for every \(i\),
contradicting the copy violation. Hence \(P\ne0\).

The homogeneous degree-\(L\) parts are equal:
\[
    \prod_{i=1}^{L}(Z+Y\ext\tau_i)
       =\prod_{i=1}^{L}(Z+Y\ext\tau_{\sigma,i}).
\]
It follows that \(P\) has total degree at most \(L-1\).
The independent uniform challenges \(\ext y,\ext z\) therefore
satisfy \(P(\ext y,\ext z)=0\) with probability at most
\(\min\{1,(L-1)/|\ExtensionField|\}\).
This joint analysis already accounts for the random mixing
of values and labels.

\noindent\textbf{Product-tree soundness.}
For one grand product, unfold the binary tree from its root
to the length-\(L\) input vector.
At a level with \(2^\ell\) parent entries, where
\(\ell\in\{0,\ldots,\mu-1\}\), let
\(\ext{\bm x}\in\ExtensionField^\ell\) be the current
evaluation point.
The weighted product identity uses the polynomial:
\[
    \left\langle
       \tensorspan{\ext{\bm x}},\tensorspan{\bm X}
    \right\rangle
    \left\langle\ext{\bm u},\tensorspan{\bm X}\right\rangle
    \left\langle\ext{\bm v},\tensorspan{\bm X}\right\rangle,
\]
where \(\ext{\bm u},\ext{\bm v}\) are the two child subvectors
and \(\bm X\) denotes \(\ell\) formal variables.
Its degree in each variable is at most three.
Sum-check reduces its Boolean sum to child evaluations
\(\ext v_L,\ext v_R\) at one point \(\ext{\bm t}\), with
error at most \(3\ell/|\ExtensionField|\).

After both values are fixed, sample a fresh
\(\ext\eta\sample\ExtensionField\) and form the next claim:
\[
    (1-\ext\eta)\ext v_L+\ext\eta\ext v_R
       =
    \left\langle
       \ext{\bm u}\mathbin{\|}\ext{\bm v},
       \tensorspan{(\ext\eta,\ext{\bm t})}
    \right\rangle.
\]
If either child claim is false, this merge hides their error
with probability at most \(1/|\ExtensionField|\).
Thus one tree has reduction error at most:
\[
    \sum_{\ell=0}^{\mu-1}
       \frac{3\ell+1}{|\ExtensionField|}
       =
    \frac{3\mu(\mu-1)/2+\mu}{|\ExtensionField|}.
\]
The \(\ell=0\) level has no sum-check rounds but still uses
a fresh merge challenge.
Each tree ends in a single MLE claim on its original input.
Subtracting the known shift \(\ext z\) and expanding the
value-label combination produces three wire claims per tree
and one permuted-label claim across the two trees.

\noindent\textbf{Openings and composition.}
The gate argument uses eight direct openings and two packed
openings. The copy argument uses seven direct openings.
Fourteen direct openings concern length-\(m\) vectors and
one concerns \(\ext{\bm\tau}_\sigma\) of length \(L\).
Each packed opening terminates in one underlying opening at
length \(m\). Consequently the total opening error is at most
\(15\epsilon_{\mathsf{PCS}}+2\epsilon_{\mathsf{open}}(N)\).
All queries may be adaptive under the stated PCS assumption.
A union bound with the gate, fingerprint, two tree, and binding
errors proves Eq.~\eqref{eq:together-multiple-soundness}.

\noindent\textbf{Efficiency.}
Computing \(\bm W_L\circ\bm W_R\) costs \(O(N)\bb\).
The one packed Hadamard proof and two packed openings cost
\(O((\kappa+k)N)\bb\), using \(m\geq\kappa\).
The gate sum-check operates on a constant number of
length-\(m\) extension-field tables and costs
\(O(m)\) extension-field operations.

For the product trees, the layer lengths form a geometric
series summing to \(O(L)\). Forming their entries, generating
the equality tables, and running the layer sum-checks therefore
cost \(O(L+(\log_2L)^2)\) extension-field operations in total.
The layer transcripts and verifier checks cost
\(O((\log_2L)^2)\) extension-field elements and operations,
respectively.
Schoolbook multiplication converts these bounds to the
\(\kappa^2[L+(\log_2L)^2]\bb\) term.

The packed Hadamard check contributes \(O(\kappa^3)\bb\)
to verifier time, and the two packed opening reductions add
\(O(\kappa^2[\kappa+k\log_2m])\bb\).
Packing the public I/O reads \(O(\kappa n_{\mathsf{io}})\)
base-field coordinates, and its terminal contractions cost
\(O(n_{\mathsf{io}}+\log L)\) extension-field operations.
The identity labels are handled by
Proposition~\ref{proposition:position-encoding-evaluation}.
Combining these costs proves the verifier bound.

The packed messages occupy \(O(\kappa^2+(\kappa+k)\log_2m)\)
base-field coordinates, in addition to the
\(O((\log_2L)^2)\) extension-field elements of the direct
gate and copy reductions.
Sequential execution stores \(O(N)\) witness coordinates and
\(O(L)\) extension-field product-tree entries, giving
\(O(N+\kappa L)\) coordinates of space.
These are the bounds in
Eq.~\eqref{eq:together-multiple-efficiency}.
The three witness commitments and seventeen underlying openings
are charged separately, at the lengths stated in the theorem.
For \(L=O(m)\), substituting \(N=m\kappa\) and dividing the
total prover cost by \(\kappa\) gives the amortized bound.
\end{proof}
